\uuid{Plsl}
\exo7id{93}
\titre{exo7 93}
\auteur{bodin}
\organisation{exo7}
\datecreate{1998-09-01}
\isIndication{false}
\isCorrection{false}
\chapitre{Nombres complexes}
\sousChapitre{Autre}
\cpge{mpsi-mp2i.nombres-complexes.b, pcsi.nombres-complexes.b, ptsi.nombres-complexes.b, mpsi-mp2i.nombres-complexes.c, pcsi.nombres-complexes.c, ptsi.nombres-complexes.c, bcpst1.nombres-complexes.a, bcpst1.nombres-complexes.b}
\module{Algèbre}
\niveau{L1}
\difficulte{}

\contenu{
\question{
Montrer que tout nombre complexe $z$ non r\'eel de module $1$ peut se mettre sous
la forme $\frac{1+ir}{1-ir}$, o\`u $r \in \Rr$.
}
}
